\section{Delooping using higher inductive types}
\label{delooping-hit}

\paragraph{Delooping as a higher inductive type}
Given a group~$G$, its delooping should have a point~$\pt{}$ and a loop for
every element of the group. Moreover, we should ensure that the multiplication
of $G$ coincides with the concatenation operation on the loop space, and that
the type we obtain is a (pointed connected) groupoid. This suggests considering
a higher inductive type, noted $\K(G,1)$, with the following constructors
\[
  \begin{array}{l@{\ :\ }l}
    \pt{}&{\K(G,1)}\\
    \cloop&G\to\pt{}=\pt{}\\
    \cloopcomp&(x,y:G)\to\cloop x\pcomp\cloop y=\cloop(xy)\\
    % \cloopid&\cloop 1=\refl{}\\
    \ctrunc&\isGroupoid(\K(G,1))
  \end{array}
\]
This construction was first proposed by Finster and Licata. They also showed,
using the encode-decode method, that it is a delooping of the original group,
\ie $\Loop\K(G,1)=G$, see~\cite[Theorem~3.2]{licata2014eilenberg}. Note that we
only ask here that $\cloop$ preserves multiplication (with $\cloopcomp$),
because it can be shown that this implies preservation of unit and inverses. In
particular, preservation of unit (see \agda{EM}{loop-id}) renders superfluous
one of the constructors present in the original
definition~\cite{licata2014eilenberg}.

In the following, we will define a variant of this higher inductive type when
the group~$G$ is presented, which is smaller and gives rise to computations
closer to traditional group theory.

\paragraph{Presentations of groups}
Given a set~$X$, we write $\freegroup{X}$ for the \emph{free group}
over~$X$~\cite[Theorem~6.11.6]{hottbook}. There is an inclusion
function~$\incl:X\to\freegroup{X}$ which, by precomposition, induces an
equivalence between morphisms of groups~$\freegroup{X}\to G$ and
functions~$X\to G$. We write $\freegroup{f}:\freegroup{X}\to G$ for the group
morphism thus induced by a function~$f:X\to G$.
% This group is such that, for any group~$G$, a function $X\to G$ extends
% uniquely as a group morphism~$\freegroup{X}\to G$.
The elements of~$\freegroup X$ can be described as formal
composites~$a_1\ldots a_n$ where each $a_i$ is an element of~$X$ or a formal
inverse of an element of~$X$ (such that an element with an adjacent formal
inverse cancel out).

Any free group~$\freegroup X$ admits a delooping as a wedge of an $X$-indexed familly of circles. The corresponding type~$\bigvee_X\S1$ can be described as the coequalizer
\begin{equation}
  \label{wedge-of-circles}
  \begin{tikzcd}
    X\ar[r,shift left]\ar[r,shift right]&1\ar[r,dotted]&\bigvee_X\S1
  \end{tikzcd}
\end{equation}
or, equivalently, as the higher inductive type
% \[
  % \begin{array}{l@{\ :\ }l}
    % \pt{}&\bigvee_X\S1\\
    % \cloop&X\to\pt{}=\pt{}
  % \end{array}
  % \]
generated by the two constructors $\pt{}:\bigvee_X\S1$ and $\cloop:X\to\pt{}=\pt{}$.

\begin{proposition}
  \label{bouquet-loop-space}
  We have $\Loop\bigvee_X\S1=\freegroup{X}$, \ie the above type is a $\B\freegroup X$.
\end{proposition}
\begin{proof}
  The fact that $\bigvee_X\S1$ is a delooping of~$\freegroup X$ is not too difficult to show when~$X$
  has decidable equality, see~\cite[Exercise 8.2]{hottbook}
  and~\cite{kraus2018free}, but the general case is more involved and was recently proved
  in~\cite{warn2023path}: the main issue is to show that this type is a groupoid.
\end{proof}

A group \emph{presentation} $\pres XR$ consists of a set~$X$ of
\emph{generators}, a set~$R$ of \emph{relations}, and two functions
$\fst,\snd:R\to\freegroup X$ respectively associating to a relation its
\emph{source} and \emph{target}. We often write $r:u\To v$ for a
relation~$r$ with $u$ as source and~$v$ as target.
%
Given such a presentation~$P$, the corresponding \emph{presented} group~$\pgroup P$ is the set quotient $\freegroup{X}/R$ of the free group on $X$ under the
smallest congruence identifying the source and the target of every
relation~$r:R$. This type can be described as the type
$\pgroup P\defd\strunc{{\freegroup X}\hq R}$ obtained by taking the set
truncation of the coequalizer
\[
  \begin{tikzcd}
    R
    \ar[r,shift left,"\fst"]
    \ar[r,shift right,"\snd"']
    &
    \freegroup X
    \ar[r,dotted,"\kappa"]
    &
    \freegroup X\hq R
  \end{tikzcd}
\]
From this also follows a description of~$\pgroup P$ as a higher inductive type:
\[
  \begin{array}{l@{\ :\ }l}
    \constructor{word}&\freegroup X\to\pgroup P\\
    \constructor{rel}&(r:R)\to\constructor{word}(\fst(r))=\constructor{word}(\snd(r))\\
    \ctrunc&\isSet(\pgroup P)
  \end{array}
\]

\paragraph{A smaller delooping}
Suppose given a group~$G$ along with a presentation $P\defd\pres XR$, \ie such that $G=\pgroup P$. We define the type $\B P$ as the following higher inductive type:
\[
  \begin{array}{l@{\ :\ }l}
    \pt{}&{\B P}\\
    \cgen&X\to(\pt{}=\pt{})\\
    \crel&(r:R)\to(\freegroup\cgen(\fst(r))=\freegroup\cgen(\snd(r)))\\
    \ctrunc&\isGroupoid(\B P)
  \end{array}
\]
This type is generated by a point~$\pt{}$, then the constructor~$\cgen$ adds a
loop~$\ul a:\pt{}=\pt{}$ for every generator~$a$, the constructor~$\crel$ adds
an equality
$\ul a_1\pcomp\ul a_2\pcomp\ldots\pcomp\ul a_n=\ul b_1\pcomp\ul
b_2\pcomp\ldots\pcomp\ul b_m$ for each relation
$a_1\ldots a_n\To b_1\ldots b_m$, and the constructor~$\ctrunc$ formally takes
the groupoid truncation of the resulting type. Note that, because of the
presence of $\freegroup\cgen:\freegroup X\to(\pt{}=\pt{})$ in the type
of~$\crel$, the above inductive type is not accepted as is in standard proof
assistants such as Agda. However, a definition can be done in two stages, by
first considering $\bigvee_X\S1$ (\ie the type generated only by $\pt{}$ and $\cgen$), and then defining a
second inductive type further quotienting this type (\ie adding the
constructors $\crel$ and $\ctrunc$), see \agda{EM}{Delooping}. Also, the
definition of $\freegroup\cgen$ requires the group structure on $\bigvee_X\S1$: the group operations are easily defined from operations on paths (reflexivity, concatenation, symmetry), but the fact that it is a groupoid is non-trivial (see \cref{bouquet-loop-space}).
%
Our main result in this section is the following:

\begin{theorem}[\agda{EM}{theorem}]
  \label{presented-delooping}
  Given a presentation~$P\defd\pres XR$, the type $\B P$ is a delooping of the group~$\pgroup P$.
\end{theorem}
\begin{proof}
  % agda -> latex
  % loop -> g'
  % path -> gen*
  %
  By induction on~$\B P$, we can define a function $f:\B P\to\K(\pgroup P,1)$ such that
  $f\pt{}\defd\pt{}$, and $f(\cgen a)\defd\cloop\pgroup a$ for all $a:X$. It can
  be shown that $f$ is then such that
  $\cong f(\freegroup\cgen u)=\cloop\pgroup u$, for any $u:\freegroup X$. We can
  therefore define the image $\cong f(r)$ on a relation $r:u\To v$ as the
  composite of equalities
  \[
    \cong f(\freegroup\cgen u)=\cloop\pgroup u=\cloop\pgroup v=\cong f(\freegroup\cgen v)
  \]
  where the equality in the middle follows from the fact that we have
  $\pgroup u=\pgroup v$ because of the relation~$r$.

  In the other direction, the group morphism $\freegroup\cgen:\freegroup{X}\to \Loop\B P$ preserves relations (by~$\crel$), and thus induces a quotient morphism $g':\pgroup P\to\Loop\B P$.
  We can thus consider the function $g:\K(\pgroup P,1)\to\B P$ such that $g(\pt{})=\pt{}$, for $x:\pgroup P$ we have $\ap g(\cloop x)=g'(x)$, and for $x,y:\pgroup P$ the image of $\cloopcomp x\,y$ is canonically induced by the fact that $g'$ preserves group multiplication.

  Since $\K(\pgroup P,1)$ is a groupoid, in order to show that $f(g(x))=x$ for every $x:\K(\pgroup P,1)$, it is enough to show that it holds for $x\defd\pt{}$, which is the case by definition of $f$ and $g$, and that this property is preserved under $\cloop x$ for $x:\pgroup P$, which follows from the fact that we have $\ap f(g'(x))=\cloop x$ for any~$x:\pgroup P$ (this is easily shown by induction on~$x$).
  %
  Conversely, we have to show that $g(f(x))=x$ holds for $x:\B P$. Again, this is shown by induction on~$x$.
  %
  % \TODO{terminer ceci}
  %
  % $\sigma$ is a section of $\strunq-\circ\kappa:\freegroup X\to\pgroup P$, and
  % the fact that $g(f(x))=x$ for $x:\B P$ follows from the fact that $\sigma$
  % preserves generators.
\end{proof}

\noindent
As an interesting remark, the careful reader will note that the fact that the types $X$ and $R$ are sets does not play a role in the proof: in fact, those assumptions can be dropped here.
%
Also, note that we do not need the choice of a representative in $\freegroup X$ for every element of $\pgroup P$ in order to define the function $g$ from $\freegroup\cgen$ in the above proof: intuitively, this is because the induced function~$g$ does not depend on such a choice of representatives.
%
Finally, we should mention here that a similar result is mentioned as an exercise in \cite[Example 8.7.17]{hottbook}; the proof suggested there is more involved since it is based on a generalized van Kampen theorem.

\begin{example}
  The dihedral group $D_5$, see \cref{Dn}, admits the presentation
  \[
    \pres{r,s}{r^4=srs,sr^2s=r^3,rsr=s,r^3s=sr^2,sr^3=r^2s,s^2=1}
  \]
  Hence, by \cref{presented-delooping} we can construct a delooping of $D_5$ as an
  higher inductive type generated by two loops (corresponding to~$r$ and $s$)
  and six $2$-dimensional cells (corresponding to the relations). Note that this
  is much smaller than $\K(D_5,1)$ (it has 2 instead of 10 generating loops, and
  6 instead of 100 relations), thus resulting in shorter proofs when reasoning
  by induction.
\end{example}

\begin{example}
  Any group~$G$ admits a presentation, the \emph{standard presentation}, with one generator $\ul a$ for every element $a:G$, and relations $\ul a\,\,\ul b=\ul{ab}$ for every pair of generators, as well as $\ul 1=1$. By applying \cref{presented-delooping}, we actually recover the inductive type~$\K(G,1)$ as delooping of~$G$.
\end{example}

% Local Variables:
% TeX-master: "article"
% End:
